% mlp_edge_diagram.tex
% Schéma hranového MLP modelu (EdgeMLP)
% Použitie: \resizebox{0.72\textwidth}{!}{% mlp_edge_diagram.tex
% Schéma hranového MLP modelu (EdgeMLP)
% Použitie: \resizebox{0.72\textwidth}{!}{% mlp_edge_diagram.tex
% Schéma hranového MLP modelu (EdgeMLP)
% Použitie: \resizebox{0.72\textwidth}{!}{% mlp_edge_diagram.tex
% Schéma hranového MLP modelu (EdgeMLP)
% Použitie: \resizebox{0.72\textwidth}{!}{\input{figures/tikz/mlp_edge_diagram}}
\begin{tikzpicture}[
    font=\small,
    box/.style={
        draw, rounded corners=5pt, align=center,
        minimum width=3.4cm, minimum height=0.82cm,
        inner sep=4pt
    },
    inputbox/.style={box, fill=gray!10, draw=gray!50},
    layerbox/.style={box, fill=blue!8, draw=blue!40},
    outbox/.style={box, fill=green!8, draw=green!50},
    labelbox/.style={font=\scriptsize\itshape, text=gray!60, align=left,
                     text width=4.0cm},
    arrow/.style={-{Latex[length=2.5mm]}, thick}
]

\node[inputbox] (phi) at (0, 0)
    {\(\phi_{ij} \in \mathbb{R}^{10}\)\\[-2pt]
     \scriptsize vektor príznakov hrany};

\node[layerbox] (h1) at (0, -1.55)
    {Lineárna vrstva\\[-2pt]
     \scriptsize + aktivácia ReLU};

\node[layerbox] (hdots) at (0, -3.10)
    {\(\cdots\)\\[-2pt]
     \scriptsize (opakuje sa \(D\)-krát)};

\node[layerbox] (hD) at (0, -4.65)
    {Lineárna vrstva\\[-2pt]
     \scriptsize + aktivácia ReLU};

\node[outbox] (out) at (0, -6.15)
    {Záverečná projekcia\\[-2pt]
     \scriptsize logit \(s_{ij} \in \mathbb{R}\)};

\draw[arrow] (phi)   -- (h1);
\draw[arrow] (h1)    -- (hdots);
\draw[arrow] (hdots) -- (hD);
\draw[arrow] (hD)    -- (out);

\node[labelbox, anchor=north west] at (2.5, -0.25)
    {\textbf{Vstup:} 10-rozmerný vektor
     obsahujúci polohu oboch vrcholov,
     vzdialenosť, relatívnu geometriu
     a uhol hrany.};

\node[labelbox, anchor=north west] at (2.5, -2.00)
    {\textbf{Skryté vrstvy:} každá vrstva
     lineárne transformuje vstup
     a aplikuje ReLU nelinearitu.};

\node[labelbox, anchor=north west] at (2.5, -5.25)
    {\textbf{Výstup:} jediné číslo
     (logit) vyjadrujúce
     odhadovanú príslušnosť
     hrany k optimálnej trase.};

\node[draw=gray!35, rounded corners=4pt, fill=gray!5,
      font=\scriptsize, align=center, text width=6.5cm,
      inner sep=5pt] at (0, -7.55)
    {Každá hrana \(\{i,j\}\) sa spracováva \textbf{nezávisle} —
     model pri hodnotení hrany nemá prístup k informácii
     o ostatných hranách ani vrcholoch inštancie.};

\end{tikzpicture}
}
\begin{tikzpicture}[
    font=\small,
    box/.style={
        draw, rounded corners=5pt, align=center,
        minimum width=3.4cm, minimum height=0.82cm,
        inner sep=4pt
    },
    inputbox/.style={box, fill=gray!10, draw=gray!50},
    layerbox/.style={box, fill=blue!8, draw=blue!40},
    outbox/.style={box, fill=green!8, draw=green!50},
    labelbox/.style={font=\scriptsize\itshape, text=gray!60, align=left,
                     text width=4.0cm},
    arrow/.style={-{Latex[length=2.5mm]}, thick}
]

\node[inputbox] (phi) at (0, 0)
    {\(\phi_{ij} \in \mathbb{R}^{10}\)\\[-2pt]
     \scriptsize vektor príznakov hrany};

\node[layerbox] (h1) at (0, -1.55)
    {Lineárna vrstva\\[-2pt]
     \scriptsize + aktivácia ReLU};

\node[layerbox] (hdots) at (0, -3.10)
    {\(\cdots\)\\[-2pt]
     \scriptsize (opakuje sa \(D\)-krát)};

\node[layerbox] (hD) at (0, -4.65)
    {Lineárna vrstva\\[-2pt]
     \scriptsize + aktivácia ReLU};

\node[outbox] (out) at (0, -6.15)
    {Záverečná projekcia\\[-2pt]
     \scriptsize logit \(s_{ij} \in \mathbb{R}\)};

\draw[arrow] (phi)   -- (h1);
\draw[arrow] (h1)    -- (hdots);
\draw[arrow] (hdots) -- (hD);
\draw[arrow] (hD)    -- (out);

\node[labelbox, anchor=north west] at (2.5, -0.25)
    {\textbf{Vstup:} 10-rozmerný vektor
     obsahujúci polohu oboch vrcholov,
     vzdialenosť, relatívnu geometriu
     a uhol hrany.};

\node[labelbox, anchor=north west] at (2.5, -2.00)
    {\textbf{Skryté vrstvy:} každá vrstva
     lineárne transformuje vstup
     a aplikuje ReLU nelinearitu.};

\node[labelbox, anchor=north west] at (2.5, -5.25)
    {\textbf{Výstup:} jediné číslo
     (logit) vyjadrujúce
     odhadovanú príslušnosť
     hrany k optimálnej trase.};

\node[draw=gray!35, rounded corners=4pt, fill=gray!5,
      font=\scriptsize, align=center, text width=6.5cm,
      inner sep=5pt] at (0, -7.55)
    {Každá hrana \(\{i,j\}\) sa spracováva \textbf{nezávisle} —
     model pri hodnotení hrany nemá prístup k informácii
     o ostatných hranách ani vrcholoch inštancie.};

\end{tikzpicture}
}
\begin{tikzpicture}[
    font=\small,
    box/.style={
        draw, rounded corners=5pt, align=center,
        minimum width=3.4cm, minimum height=0.82cm,
        inner sep=4pt
    },
    inputbox/.style={box, fill=gray!10, draw=gray!50},
    layerbox/.style={box, fill=blue!8, draw=blue!40},
    outbox/.style={box, fill=green!8, draw=green!50},
    labelbox/.style={font=\scriptsize\itshape, text=gray!60, align=left,
                     text width=4.0cm},
    arrow/.style={-{Latex[length=2.5mm]}, thick}
]

\node[inputbox] (phi) at (0, 0)
    {\(\phi_{ij} \in \mathbb{R}^{10}\)\\[-2pt]
     \scriptsize vektor príznakov hrany};

\node[layerbox] (h1) at (0, -1.55)
    {Lineárna vrstva\\[-2pt]
     \scriptsize + aktivácia ReLU};

\node[layerbox] (hdots) at (0, -3.10)
    {\(\cdots\)\\[-2pt]
     \scriptsize (opakuje sa \(D\)-krát)};

\node[layerbox] (hD) at (0, -4.65)
    {Lineárna vrstva\\[-2pt]
     \scriptsize + aktivácia ReLU};

\node[outbox] (out) at (0, -6.15)
    {Záverečná projekcia\\[-2pt]
     \scriptsize logit \(s_{ij} \in \mathbb{R}\)};

\draw[arrow] (phi)   -- (h1);
\draw[arrow] (h1)    -- (hdots);
\draw[arrow] (hdots) -- (hD);
\draw[arrow] (hD)    -- (out);

\node[labelbox, anchor=north west] at (2.5, -0.25)
    {\textbf{Vstup:} 10-rozmerný vektor
     obsahujúci polohu oboch vrcholov,
     vzdialenosť, relatívnu geometriu
     a uhol hrany.};

\node[labelbox, anchor=north west] at (2.5, -2.00)
    {\textbf{Skryté vrstvy:} každá vrstva
     lineárne transformuje vstup
     a aplikuje ReLU nelinearitu.};

\node[labelbox, anchor=north west] at (2.5, -5.25)
    {\textbf{Výstup:} jediné číslo
     (logit) vyjadrujúce
     odhadovanú príslušnosť
     hrany k optimálnej trase.};

\node[draw=gray!35, rounded corners=4pt, fill=gray!5,
      font=\scriptsize, align=center, text width=6.5cm,
      inner sep=5pt] at (0, -7.55)
    {Každá hrana \(\{i,j\}\) sa spracováva \textbf{nezávisle} —
     model pri hodnotení hrany nemá prístup k informácii
     o ostatných hranách ani vrcholoch inštancie.};

\end{tikzpicture}
}
\begin{tikzpicture}[
    font=\small,
    box/.style={
        draw, rounded corners=5pt, align=center,
        minimum width=3.4cm, minimum height=0.82cm,
        inner sep=4pt
    },
    inputbox/.style={box, fill=gray!10, draw=gray!50},
    layerbox/.style={box, fill=blue!8, draw=blue!40},
    outbox/.style={box, fill=green!8, draw=green!50},
    labelbox/.style={font=\scriptsize\itshape, text=gray!60, align=left,
                     text width=4.0cm},
    arrow/.style={-{Latex[length=2.5mm]}, thick}
]

\node[inputbox] (phi) at (0, 0)
    {\(\phi_{ij} \in \mathbb{R}^{10}\)\\[-2pt]
     \scriptsize vektor príznakov hrany};

\node[layerbox] (h1) at (0, -1.55)
    {Lineárna vrstva\\[-2pt]
     \scriptsize + aktivácia ReLU};

\node[layerbox] (hdots) at (0, -3.10)
    {\(\cdots\)\\[-2pt]
     \scriptsize (opakuje sa \(D\)-krát)};

\node[layerbox] (hD) at (0, -4.65)
    {Lineárna vrstva\\[-2pt]
     \scriptsize + aktivácia ReLU};

\node[outbox] (out) at (0, -6.15)
    {Záverečná projekcia\\[-2pt]
     \scriptsize logit \(s_{ij} \in \mathbb{R}\)};

\draw[arrow] (phi)   -- (h1);
\draw[arrow] (h1)    -- (hdots);
\draw[arrow] (hdots) -- (hD);
\draw[arrow] (hD)    -- (out);

\node[labelbox, anchor=north west] at (2.5, -0.25)
    {\textbf{Vstup:} 10-rozmerný vektor
     obsahujúci polohu oboch vrcholov,
     vzdialenosť, relatívnu geometriu
     a uhol hrany.};

\node[labelbox, anchor=north west] at (2.5, -2.00)
    {\textbf{Skryté vrstvy:} každá vrstva
     lineárne transformuje vstup
     a aplikuje ReLU nelinearitu.};

\node[labelbox, anchor=north west] at (2.5, -5.25)
    {\textbf{Výstup:} jediné číslo
     (logit) vyjadrujúce
     odhadovanú príslušnosť
     hrany k optimálnej trase.};

\node[draw=gray!35, rounded corners=4pt, fill=gray!5,
      font=\scriptsize, align=center, text width=6.5cm,
      inner sep=5pt] at (0, -7.55)
    {Každá hrana \(\{i,j\}\) sa spracováva \textbf{nezávisle} —
     model pri hodnotení hrany nemá prístup k informácii
     o ostatných hranách ani vrcholoch inštancie.};

\end{tikzpicture}
